\begin{pastexam}[label={px:ntu-110-ie-ds-20}]{臺大工工所 110 學年度 第 20 題}
（複選；每題 5 分，每答錯一個選項倒扣 1 分至該題 0 分止）

Consider the red-black tree in Figure 2, which of the following statements are true?\\
(Assume bottom-up insertion/deletion)
\begin{enumerate}[label=(\Alph*)]
\item In Figure 2, after inserting key 14, then there are 6 red nodes.
\item In Figure 2, after inserting key 54, then key 49 and key 54 are siblings.
\item In Figure 2, after deleting key 22, then key 25 and key 73 are siblings.
\item In Figure 2, after deleting key 66, then key 35 is red.
\item In Figure 2, after deleting key 83, then key 89 is red.
\end{enumerate}
\begin{center}
\begin{tikzpicture}[b/.style={circle,draw,thick,inner sep=1pt,minimum size=6mm,font=\small\sffamily},
  r/.style={b,draw=red,dotted,very thick},>={Stealth[length=2mm]},
  eb/.style={thick,->},er/.style={very thick,red,dotted,->}]
\node[b] (58) at (0,0) {58};
\node[r] (35) at (-2.4,-1.1) {35}; \node[r] (73) at (2.4,-1.1) {73};
\node[b] (22) at (-3.45,-2.5) {22}; \node[b] (49) at (-1.35,-2.5) {49};
\node[b] (66) at (1.3,-2.5) {66}; \node[b] (83) at (3.5,-2.5) {83};
\node[r] (25) at (-2.85,-3.7) {25}; \node[r] (50) at (-0.75,-3.7) {50}; \node[r] (89) at (4.25,-3.7) {89};
\foreach \a/\b in {58/35,58/73,22/25,49/50,83/89} \draw[er] (\a) -- (\b);
\foreach \a/\b in {35/22,35/49,73/66,73/83} \draw[eb] (\a) -- (\b);
\end{tikzpicture}
\\*[4pt] Figure 2: Red-black tree
\end{center}
\end{pastexam}

\begin{pastsoln}
答案是 (A)、(B)、(D)。

圖中虛線紅框是紅色節點：$35$、$73$、$25$、$50$、$89$ 是紅的，其他是黑的。 先確認它是合法的紅黑樹：紅色節點的子節點都是黑的
（或 NIL），每條從根到 NIL 的路徑都經過 $2$ 個黑節點（NIL 不算；例如 $58 \to 35 \to 22 \to$ NIL 經過 $58, 22$）。
每個選項都從 Figure 2 的原樹開始，插入、刪除用資結書第~10.3.4、10.3.7~節的由下而上做法（虛擬碼~10.6、10.7、10.8）；
刪除有兩個子節點的節點時改刪後繼，但本題刪的節點最多只有一個子節點，用不到。 下圖中紅色節點同樣畫成虛線紅框。

\begin{center}
\begin{tikzpicture}[level distance=0.85cm,level 1/.style={sibling distance=2.4cm},level 2/.style={sibling distance=1.2cm},
  level 3/.style={sibling distance=0.7cm},
  b/.style={circle,draw=black,solid,thick,inner sep=0pt,minimum size=5.5mm,font=\footnotesize\sffamily},
  r/.style={circle,draw=red,dotted,very thick,inner sep=0pt,minimum size=5.5mm,font=\footnotesize\sffamily},
  eb/.style={draw=black,solid,thick},er/.style={draw=red,dotted,very thick}]
\node[b] {58} child {node[r] {35} child {node[b] {22} child {node[r] {14} child[missing] child[missing] edge from parent[er]} child {node[r] {25} child[missing] child[missing] edge from parent[er]} edge from parent[eb]} child {node[b] {49} child[missing] child {node[r] {50} child[missing] child[missing] edge from parent[er]} edge from parent[eb]} edge from parent[er]} child {node[r] {73} child {node[b] {66} child[missing] child[missing] edge from parent[eb]} child {node[b] {83} child[missing] child {node[r] {89} child[missing] child[missing] edge from parent[er]} edge from parent[eb]} edge from parent[er]};
\node[font=\small] at (0,-3.2) {(A) 插入 14};
\end{tikzpicture}\hspace{0.4cm}\begin{tikzpicture}[level distance=0.85cm,level 1/.style={sibling distance=2.4cm},level 2/.style={sibling distance=1.2cm},
  level 3/.style={sibling distance=0.7cm},
  b/.style={circle,draw=black,solid,thick,inner sep=0pt,minimum size=5.5mm,font=\footnotesize\sffamily},
  r/.style={circle,draw=red,dotted,very thick,inner sep=0pt,minimum size=5.5mm,font=\footnotesize\sffamily},
  eb/.style={draw=black,solid,thick},er/.style={draw=red,dotted,very thick}]
\node[b] {58} child {node[r] {35} child {node[b] {22} child[missing] child {node[r] {25} child[missing] child[missing] edge from parent[er]} edge from parent[eb]} child {node[b] {50} child {node[r] {49} child[missing] child[missing] edge from parent[er]} child {node[r] {54} child[missing] child[missing] edge from parent[er]} edge from parent[eb]} edge from parent[er]} child {node[r] {73} child {node[b] {66} child[missing] child[missing] edge from parent[eb]} child {node[b] {83} child[missing] child {node[r] {89} child[missing] child[missing] edge from parent[er]} edge from parent[eb]} edge from parent[er]};
\node[font=\small] at (0,-3.2) {(B) 插入 54};
\end{tikzpicture}
\end{center}
\begin{itemize}
\item (A) 對。 $14$ 沿 $58 \to 35 \to 22$ 走，放在 $22$ 的左邊，塗紅。 父節點 $22$ 是黑的，不必修正。 紅色節點變成原來的 $5$ 個加上
$14$，共 $6$ 個。
\item (B) 對。 $54$ 沿 $58 \to 35 \to 49 \to 50$ 走，放在 $50$ 的右邊，塗紅。 父節點 $50$ 是紅的，要修正：祖父 $49$，叔叔是 $49$ 的左邊
（NIL，黑的）。 $54$ 是右子節點、$50$ 也是右子節點，是情況 3 的鏡像：$50$ 改黑、$49$ 改紅，對 $49$ 左旋。 $50$ 成為 $35$ 的右子節點，
左右是 $49$、$54$（都是紅的），$49$ 和 $54$ 是兄弟。
\end{itemize}
\begin{center}
\begin{tikzpicture}[level distance=0.85cm,level 1/.style={sibling distance=2.4cm},level 2/.style={sibling distance=1.2cm},
  level 3/.style={sibling distance=0.7cm},
  b/.style={circle,draw=black,solid,thick,inner sep=0pt,minimum size=5.5mm,font=\footnotesize\sffamily},
  r/.style={circle,draw=red,dotted,very thick,inner sep=0pt,minimum size=5.5mm,font=\footnotesize\sffamily},
  eb/.style={draw=black,solid,thick},er/.style={draw=red,dotted,very thick}]
\node[b] {58} child {node[r] {35} child {node[b] {25} child[missing] child[missing] edge from parent[eb]} child {node[b] {49} child[missing] child {node[r] {50} child[missing] child[missing] edge from parent[er]} edge from parent[eb]} edge from parent[er]} child {node[r] {73} child {node[b] {66} child[missing] child[missing] edge from parent[eb]} child {node[b] {83} child[missing] child {node[r] {89} child[missing] child[missing] edge from parent[er]} edge from parent[eb]} edge from parent[er]};
\node[font=\small] at (0,-3.2) {(C) 刪除 22};
\end{tikzpicture}\hspace{0.3cm}\begin{tikzpicture}[level distance=0.85cm,level 1/.style={sibling distance=2.4cm},level 2/.style={sibling distance=1.2cm},
  level 3/.style={sibling distance=0.7cm},
  b/.style={circle,draw=black,solid,thick,inner sep=0pt,minimum size=5.5mm,font=\footnotesize\sffamily},
  r/.style={circle,draw=red,dotted,very thick,inner sep=0pt,minimum size=5.5mm,font=\footnotesize\sffamily},
  eb/.style={draw=black,solid,thick},er/.style={draw=red,dotted,very thick}]
\node[b] {58} child {node[r] {35} child {node[b] {22} child[missing] child {node[r] {25} child[missing] child[missing] edge from parent[er]} edge from parent[eb]} child {node[b] {49} child[missing] child {node[r] {50} child[missing] child[missing] edge from parent[er]} edge from parent[eb]} edge from parent[er]} child {node[r] {83} child {node[b] {73} child[missing] child[missing] edge from parent[eb]} child {node[b] {89} child[missing] child[missing] edge from parent[eb]} edge from parent[er]};
\node[font=\small] at (0,-3.2) {(D) 刪除 66};
\end{tikzpicture}\hspace{0.3cm}\begin{tikzpicture}[level distance=0.85cm,level 1/.style={sibling distance=2.4cm},level 2/.style={sibling distance=1.2cm},
  level 3/.style={sibling distance=0.7cm},
  b/.style={circle,draw=black,solid,thick,inner sep=0pt,minimum size=5.5mm,font=\footnotesize\sffamily},
  r/.style={circle,draw=red,dotted,very thick,inner sep=0pt,minimum size=5.5mm,font=\footnotesize\sffamily},
  eb/.style={draw=black,solid,thick},er/.style={draw=red,dotted,very thick}]
\node[b] {58} child {node[r] {35} child {node[b] {22} child[missing] child {node[r] {25} child[missing] child[missing] edge from parent[er]} edge from parent[eb]} child {node[b] {49} child[missing] child {node[r] {50} child[missing] child[missing] edge from parent[er]} edge from parent[eb]} edge from parent[er]} child {node[r] {73} child {node[b] {66} child[missing] child[missing] edge from parent[eb]} child {node[b] {89} child[missing] child[missing] edge from parent[eb]} edge from parent[er]};
\node[font=\small] at (0,-3.2) {(E) 刪除 83};
\end{tikzpicture}
\end{center}
\begin{itemize}
\item (C) 錯。 $22$ 是黑的，只有一個子節點 $25$（紅）。 拿掉 $22$，$25$ 接到它的位置；被拿掉的是黑節點、接上來的 $x = 25$ 是紅的，
把 $25$ 改黑就恢復黑高度。 $25$ 是 $35$ 的左子節點，兄弟是 $49$；$73$ 是 $35$ 的兄弟，深度不同，不是 $25$ 的兄弟。
\item (D) 對。 $66$ 是黑色的葉節點，拿掉後 $x = $ NIL（雙重黑），$p = 73$，$x$ 是左子節點，兄弟 $s = 83$ 是黑的，$s$ 的右子節點 $89$
是紅的：情況 4。 $s$ 改成 $p$ 的顏色（紅），$p = 73$ 改黑，$89$ 改黑，對 $73$ 左旋。 $83$ 成為 $58$ 的右子節點（紅），左右是 $73$、$89$
（黑）。 整個過程只動了 $58$ 的右子樹，$35$ 仍然是紅的。
\item (E) 錯。 $83$ 是黑的，只有一個子節點 $89$（紅）。 拿掉 $83$，$89$ 接上來；被拿掉的是黑節點，$x = 89$ 是紅的，改黑即可。 所以
$89$ 變成黑的。
\end{itemize}
\end{pastsoln}
